% =====================================================================
%  chunks/ch04g.tex —— 原书 p.216–226（PDF 223–233）
%  第 4 章 §3（下）：球面同伦群阶的估计、Milnor 例子、配边的整性
% =====================================================================
% [原文 p.216]
由此可知，$G$-流形的等价类与同伦群 $\pi_{m+1}(M(G))$ 的元素之间成立自然的一一对应。

若把注意力限制在法丛具有指定结构群的流形 $M^k\subset S^{m+k}$ 上，就得到与经典李群相对应的下列配边环：

\begin{itemize}[label={},leftmargin=2.4em,itemsep=2pt]
\item $\Omega_*^{O},\ \Omega_*^{SO}$（经典配边环，20 世纪 50 年代初引入）；
\item $\Omega_*^{fr},\ \Omega_q^{fr}\cong\pi_{q+n}(S^n)$；$G=\{1\}$（标架流形的配边，20 世纪 30 年代末为研究球面同伦群而引入）；
\item $\Omega_*^{U},\ \Omega_*^{SU}$；$U_k,\ SU_k\subset SO_{2k}$（酉配边群与特殊酉配边群）；
\item $\Omega_q^{Sp}$；$Sp_k\subset SO_{4k}$（辛配边群）；
\item $\Omega_q^{Spin}$；$\rho:\Spin_k\to SO_k$（旋量配边群）。
\end{itemize}

% [Tutor Note: 第二行同时给出环与群：$\Omega_*^{fr}$ 是标架配边环，而 $\Omega_q^{fr}\cong\pi_{q+n}(S^n)$（$n$ 充分大）说明第 $q$ 个标架配边群就是第 $q$ 个稳定球面同伦群；这正是 Pontryagin--Thom 构造的内容，也是下文用几何方法估计球面同伦群的入口。]

对 $G=U,\,SU,\,Sp,\,Spin$ 及其进一步推广的 $G$-配边环 $\Omega_*^{G}$ 的结构研究，始于 1960 年前后（由 Milnor 与 Novikov 开创）。对经典群 $G$，流形稳定法丛的 $G$-结构由切丛的 $G$-结构决定；因此，例如对 $G=U$，一个实维数为 $2k$ 的复（或拟复）流形 $M^{2k}$ 决定了 $\Omega_{2k}^{U}$ 中的一个配边类：

\[ [M^{2k}]\in\Omega_{2k}^{U}. \]

一般的 $n+1$ 次多项式方程（在 $n+1$ 个齐次复坐标中）在 $\CP^n$ 中定义出一个 $SU$-流形 $M^{2n-2}$，它代表配边类

\[ [M^{2n-2}]\in\Omega_{2n-2}^{SU}. \]

当 $n=3$ 时，这样得到 4 实维的 Kummer 曲面 $K^4$：它几乎可平行化（即去掉一点后可平行化），单连通，并且具有下列不变量：

\begin{align}
\chi(K^4) &= c_2(K^4) = 24, \nonumber\\
\tau(K^4) &= 16, \quad p_1(K^4) = 48, \quad c_1(K^4) = 0. \tag{3.16}
\end{align}

% [原文 p.217]
如 §1 中给出 $G$-丛的示性类定义时所指出的，每个这样的示性类都由 $H^*(BG)$ 的一个元素 $\psi$ 唯一决定，其中 $BG$ 是万有 $G$-丛的底空间。于是对闭流形 $M^k$，$M^k$ 的示性数（上文已提及）就是 $k$ 维示性类在同调基本类 $[M^k]$ 上取的值 $\psi([M^k])$（$\psi\in H^{k}(BG)$）。对 $G\subset SO_m$，有（如上文所述）Thom 同构

\begin{align}
\phi &: H^{k}(BG;\Q)\cong H^{k+m}(M(G_m);\Q), \nonumber\\
\phi(z) &= z\cdot u_m .
\end{align}

% [Tutor Note: 这里 $u_m$ 是 Thom 类，$\phi$ 即「乘以 Thom 类」；示性数就是 $\psi$ 与基本类的配对，注意 $\psi$ 的次数 $k$ 与 $u_m$ 的次数 $m$ 相加后落在 $M(G_m)$ 的 $k+m$ 维上，写批注时勿把 $u_m$ 误当成 $M$ 的基本类。]

另一方面，由 Cartan--Serre 理论（见第 3 章 §7），当 $k<m-1$ 时 Hurewicz 同态给出同构

\[ \pi_{m+k}(M(G))\otimes\Q\cong H^{k+m}(M(G);\Q). \]

由这两个同构可以断定，示性数给出了同伦群

\[ \pi_{m+k}(M(G))\otimes\Q \]

上的一族完全的线性形式。只要 $k<m-1$，对任何 $G\subset SO_m$，示性数都可以看作 $G$-配边类的不变量。由此可知，$\Omega_*^{G}$ 中的 $G$-配边类由示性数决定到挠为止。

作为 Thom 定理陈述 (ii) 的推论，可以证明：对 $G=SO_m$，任给闭链 $z\in H_{n-m}(M^n;\Z)$，都存在整数 $\lambda\ne 0$，使元素 $\lambda z$ 可由子流形代表（这一点其实在非稳定维数下也成立）。

对最简单的（即维数较小的）李群，Thom 空间 $M(G)$ 如下：

\begin{gather*}
G=SO_1=\{1\},\quad M(SO_1)=S^1;\\
G=O_1=\Z/2,\quad M(O_1)=\RP^\infty;\\
G=SO_2=S^1=U_1,\quad M(SO_2)=M(U_1)=\CP^\infty.
\end{gather*}

由此看出，$M(SO_1)$、$M(O_1)$ 与 $M(SO_2)=M(U_1)$ 分别是 Eilenberg--MacLane 空间 $K(\Z,1)$、$K(\Z/2,1)$ 与 $K(\Z,2)$。于是 $H_{n-1}(M^n;\Z)$、$H_{n-1}(M^n;\Z/2)$ 与 $H_{n-2}(M^n;\Z)$ 中的所有闭链都可由子流形实现。

环 $\Omega_*^{SO}\otimes\Q$ 以配边类 $[\CP^{2i}]$ 为生成元（当 $j\ne 4k$ 时群 $\Omega_j^{SO}\otimes\Q$ 为平凡）。作为 $\Omega_{4k}^{SO}\otimes\Q$ 的线性基，可以取形状为

\[ [\CP^{2i_1}]\times\cdots\times[\CP^{2i_s}],\quad i_1+\cdots+i_s=k \]

的一族配边类。Pontryagin 数给出环 $\Omega_*^{SO}\otimes\Q$ 的一族完全的不变量。符号差 $\tau(M^{4k})$ 也是配边不变量；事实上 $\tau$ 是环同态 $\Omega_*^{SO}\to\Z$，在生成元 $[\CP^{2i}]$ 上取值 1：

% [原文 p.218]
\begin{align}
\tau(M_1^{4k})\cdot\tau(M_2^{4l}) &= \tau(M_1^{4k}\times M_2^{4l}), \nonumber\\
\tau(\CP^{2i}) &= 1. \tag{3.17}
\end{align}

% [Tutor Note: (3.17) 说明 $\tau$ 是乘性的（环同态），故只须在生成元上算；而 $\tau(\CP^{2i})=1$ 与下文 $L$ 级数的常数项 1 是同一件事的两面。]

符号差 $\tau(M^{4k})$ 如何用 $M^{4k}$ 的 Pontryagin 数表示？当 $k=1$，即在群 $\Omega_4^{SO}$ 上，有上文提到的 Thom--Rohlin 公式 $\tau=\frac{1}{3}p_1$。由上文 §1 阐述示性类时给出的公式，可以得到流形 $\CP^4$ 与 $\CP^2\times\CP^2$ 的 Pontryagin 类的下列公式：

对 $\CP^4$：
\[ p=1+p_1+p_2=(1+u^2)^5, \]

其中 $u$ 是 $H^2(\CP^4;\Z)\cong\Z$ 的生成元；于是

\[ p_1=5u^2,\quad p_2=10u^4,\quad p_1^2=25u^4. \]

对 $\CP_1^2\times\CP_2^2$：
\[ p=1+p_1+p_2=(1+u_1^2)^3(1+u_2^2)^3, \]

其中 $u_i$ 生成 $H^2(\CP_i^2;\Z)$；于是

\[ p_1^2=18u_1^2u_2^2,\quad p_2=9u_1^2u_2^2. \]

这些流形的 Pontryagin 数（即上述类在基本类上的取值）见下表：

\begin{center}
\begin{tabular}{lccc}
\toprule
 & $\tau$ & $p_1^2$ & $p_2$ \\
\midrule
$\CP^4$ & 1 & 25 & 10 \\
$\CP^2\times\CP^2$ & 1 & 18 & 9 \\
\bottomrule
\end{tabular}
\end{center}

于是推出符号差 $\tau$ 用 Pontryagin 数表示的下列表达式（在群 $\Omega_8^{SO}$ 上）：

\begin{equation}
\tau=\frac{1}{45}(7p_2-p_1^2). \tag{3.18}
\end{equation}

Hirzebruch 于 20 世纪 50 年代中期给出的这类一般公式如下：

\begin{equation}
\tau(M^{4k})=\langle L_k(p_1,\dots,p_k),[M^{4k}]\rangle, \tag{3.19}
\end{equation}

其中多项式 $L_k(p_1,\dots,p_k)$ 由级数

\[ L=1+L_1+L_2+\cdots+L_k+\cdots=\prod_{i=1}^{N}\frac{u_i}{\tanh u_i},\qquad N\gg k \]

的 $k$ 次齐次部分给出，而该级数被表为前 $k$ 个初等对称函数的多项式：

% [原文 p.219]
\[ p_j=\sum_{i_1<\cdots<i_j}u_{i_1}^2\cdots u_{i_j}^2. \]

这是 §1 讨论示性类时提到的通用方法的一个特例：用一个形式幂级数

\[ Q(u)=1+Q_1u^2+Q_2u^4+\cdots \]

得到 Pontryagin 示性类中的乘性表达式。对上面的 Hirzebruch 公式，取形式级数

\begin{equation}
L(u^2)=\frac{u}{\tanh u}. \tag{3.20}
\end{equation}

任何由 Pontryagin 数（在 $\Omega_*^{SO}\otimes\Q$ 上）或由 Chern 数（在 $\Omega_*^{U}$ 上）给出的乘性示性类 $Q$，都由它在 $\CP^m$ 上的取值完全决定，而这些取值正给出级数 $Q(u)$。若取形式级数（乘性示性类 $Q$ 的「通有级数」）

\[ g_Q(u)=\sum_{n\ge 0}Q(\CP^n)\frac{u^{n+1}}{n+1}, \]

其形式逆为 $g_Q^{-1}(u)$，即 $g_Q^{-1}(g_Q(u))\equiv u$，则可以建立下列一般公式（属于 Novikov，20 世纪 60 年代末；与形式群理论有关——见第 3 章 §9）：

\begin{equation}
\frac{u}{Q(u)}=g^{-1}(u). \tag{3.21}
\end{equation}

% [Tutor Note: 这里 $g_Q$ 与 $g^{-1}$ 是互逆的形式幂级数，而 (3.21) 把「乘性示性类」与「形式群的逆」联系起来：$u/Q(u)$ 恰是形式群的不变微分方向的对数导数。初读易把 $g^{-1}(u)$ 误读成 $1/g(u)$，务必按复合逆理解。]

在符号差 $\tau$ 的情形，Hirzebruch 多项式 $L_k(p_1,\dots,p_k)$ 由复射影空间的符号差决定：

\[ \tau(\CP^{2k})=1,\quad \tau(\CP^{2k+1})=0. \]

现在我们再定义一个重要的不变量——Todd 亏格 $T$，它由在 $\CP^k$ 上的值

\[ T(\CP^k)=1,\quad j\ge 0, \]

以及通有级数 $g_T(u)$ 确定：

\begin{align}
g_T(u) &= \sum_{n\ge 0}\frac{u^{n+1}}{n+1}=-\log(1-u), \nonumber\\
g_T^{-1} &= 1-\exp(-u), \tag{3.22}\\
T(u) &= \frac{u}{1-\ee^{-u}},\qquad T(M^n)=\Big\langle\prod_j\frac{u_j}{1-\ee^{-u_j}},[M^n]\Big\rangle. \nonumber
\end{align}

% [Tutor Note: 原书 p.219 印作 $T(\CP^k)=1,\ j\ge 0$，其中 $j$ 无定义；按上下文应为 $k\ge 0$。正文照录原书，读者按 $k$ 理解即可。]

% [原文 p.220]
利用复代数簇与全纯纤维丛的某些一般范畴性质，可以从 Hirzebruch 符号差公式 (3.19) 推出 Riemann--Roch--Hirzebruch 定理，做法如下。设 $\xi$ 是复代数簇 $M^n$（复维数为 $n$）上的全纯向量丛，以 $\cF_\xi$ 记 $\xi$ 的全纯截面芽层。定义不变量 $\chi$ 为

\[ \chi(M^n,\xi)=\sum_{j\ge 0}(-1)^j\dim H^j(M^n;\cF_\xi). \]

可以证明，这个不变量 $\chi(M^n,\xi)$ 由下列 Riemann--Roch--Hirzebruch 公式给出：

\begin{equation}
\chi(M^n,\xi)=\langle(\ch\xi)\cdot T(M^n),[M^n]\rangle, \tag{3.23}
\end{equation}

其中 $\ch\xi$ 是 $\xi$ 的 Chern 特征，$T(M^n)$ 如 (3.22) 所定义，$n=2k$。当向量丛 $\xi$ 是一维且平凡时，得到 Todd 亏格，它也称为簇 $M^n$ 的「全纯 Euler 示性数」。在经典情形（复曲线 $M^1$ 的情形），全纯向量丛 $\xi$ 由除子 $D=\sum n_ix_i$ 的等价类决定，其中 $x_i\in M^1$，$n_i$ 是整数；此时层 $\cF_\xi$ 改记为 $\cF_D$。于是群 $H^0(M^1;\cF_D)$ 由在除子 $D$ 之外没有极点的函数组成，这些函数满足

\[ (f)+D\ge 0, \]

其中 $(f)=\sum m_ky_k$（$m_k\le 0$）是函数 $f$ 的极点除子。经典 Riemann--Roch 定理断言

\begin{gather*}
\dim H^0(M^1;\cF_D)-\dim H^1(M^1;\cF_D)=n(D)-g+1,\\
n=\sum n_j,\qquad g=\text{曲线 }M^1\text{ 的亏格}.
\end{gather*}

利用等式

\[ \dim H^1(M^1;\cF_D)=\dim H^1(M^1;\cF_{K-D}), \]

可以把这一定理写成另一种形式，其中 $K$ 是亚纯形式的零点与极点的除子。

Todd 亏格（见上文 (3.22)）在整个配边群 $\Omega_{2k}^{U}$ 上取整数值；其原因将在下文说明。20 世纪 50 年代末，Hirzebruch 与 Atiyah 证明了若干「整性定理」，它们在许多计算中很有用。其中最重要的一条如下：由 Todd 亏格的上述公式 $T=\dfrac{u}{1-\ee^{-u}}$ 有

\[ T(u)=\ee^{u/2}\frac{u/2}{\sinh u/2}. \]

% [原文 p.221]
令

\begin{equation}
A(u^2)=\frac{u/2}{\sinh u/2}=\ee^{-u/2}T(u). \tag{3.24}
\end{equation}

由这个级数按上述方式决定的示性类

\begin{align}
A(p_1,\dots,p_k,\dots) &= 1+A_1+\cdots+A_k+\cdots, \nonumber\\
A_1 &= -\frac{p_1}{24},\qquad A_2=\frac{1}{45\cdot 2^7}\bigl(-4p_2+7p_1^2\bigr), \tag{3.25}
\end{align}

称为 $A$-亏格。可以证明：只要 $w_1(M^{4k})=0$、$w_2(M^{4k})=0$，数 $\langle A_k,[M^{4k}]\rangle$ 对所有的 $k$ 都是整数，而且当 $k$ 为奇数时是偶数。由此可知，对所有的 $k$，$A_k$ 定义了旋量配边群 $\Omega_{4k}^{Spin}$ 上的一个整值线性形式。（注意：流形具有旋量结构当且仅当它可定向（$w_1=0$）并且 $w_2=0$。）

% [Tutor Note: 原书对同一个亏格混用两种记法：本节（原书 p.221）写作 $A$-亏格、$A(u^2)$、$A_k$，而第 4 章 §4（原书 p.233）与第 3 章 (8.24) 写作 $\hat A$。本书一切照原书排印，读者须知二者均指同一个亏格，即通行的 $\hat A$-亏格（Â-亏格）。]
% [Tutor Note: 原书此处先说 Todd 亏格在 $\Omega_{2k}^{U}$ 上整值（只用 $c_1$ 的整性），而 $A$-亏格只在 $w_1=w_2=0$ 的 $\Spin$ 配边群上整值；两个整性定理的适用范围不同，是本节最容易混淆之处。]

为了给出这些整性定理的进一步应用，我们需要回忆同伦群 $\pi_{j-1}(SO)$、$\pi_{j-1}(U)$ 的 Bott 周期性的一个最重要的推论，这两个群分别分类球面 $S^j$ 上的定向丛与复丛。群 $\pi_{4k}(SO)$ 的生成元对应于 $S^{4k}$ 上一个「基本」定向向量丛 $\eta^{4k}$，其 Chern 特征（见第 3 章 §8）为 $\ch\eta^{4k}=a_k\mu$，其中 $\langle\mu,[S^{4k}]\rangle=1$，且

\[ a_k=\begin{cases}1, & k\text{ 为偶数且 }k\ge 2,\\ 2, & k\text{ 为奇数}.\end{cases} \]

相应的「基本」$U$-丛 $\eta_{\C}$ 在 $S^{4k}$ 上的 Chern 特征为 $\ch\eta_{\C}=\mu$。$\eta_{\C}$ 的 Chern 类与 Pontryagin 类如下：

\begin{align}
p_k(\eta_{\C}) &= a_k\cdot(2k-1)!\mu, \nonumber\\
c_k(\eta_{\C}) &= a_k\cdot(k-1)!\mu. \tag{3.26}
\end{align}

（注意：公式 (3.26) 可由 Chern 特征的归纳定义

\[ c_k=(k-1)!\,\ch_k+f(c_1,\dots,c_{k-1}) \]

推出。）

\begin{remark}{}
若 $M^n$ 是几乎可平行化的流形，稳定法丛为 $\nu_M$，则不难证明存在一个度为 1 的「法向」映射 $f:M^n\to S^n$（即把基本类 $[M^n]$ 映为 $[S^n]$），使得 $f^*(\xi)\cong\nu_M$，其中 $\xi$ 是球面 $S^n$ 上某个稳定向量丛。这一构造表明，几乎可平行化流形的 Pontryagin 类（以及所有其他示性类）都可用球面上向量丛的示性类表示。我们还注意到，类 $p_k(M^{4k})$ 是唯一可能非平凡的 Pontryagin 类。\hfill$\square$
\end{remark}

% [原文 p.222]
由这一注可知，任何几乎可平行化流形 $M^{4k}$ 的稳定法丛（或切丛）由球面 $S^{4k}$ 上「基本」丛 $\eta$ 的某个倍数决定。特别地，唯一可能非平凡的 Pontryagin 类 $p_k(M^{4k})$ 具有形状

\[ p_k(M^{4k})=\lambda_M\cdot a_k\cdot\mu\cdot(2k-1)!,\qquad \langle\mu,[M^{4k}]\rangle=1, \]

其中 $\lambda_M$ 是某个整数。以 $\lambda_{\min}$ 记遍取所有几乎可平行化流形 $M^{4k}$ 时 $\lambda_M$ 的最小值。这个数 $\lambda_{\min}$ 决定了前面提到的群 $J\pi_{4k-1}(SO)\subset\pi_{4k-1}(S^n)$（$n>4k$）的阶，该群可通过标准球面 $S^{4k-1}\subset\R^{n+4k-1}$ 的标架来实现：

\begin{equation}
\lambda_{\min}=|J\pi_{4k-1}(SO)|. \tag{3.27}
\end{equation}

% [Tutor Note: 原书 p.222 把 $J$ 同态的像写作 $J\pi_{4k-1}(SO)\subset\pi_{4k-1}(S^n)$（$n>4k$），但当 $n>4k$ 时 $\pi_{4k-1}(S^n)=0$，与此处断言矛盾；按其本意应为稳定球面同伦群 $\pi_{n+4k-1}(S^n)$，与同页下文「$S^{4k-1}\subset\R^{n+4k-1}$ 的标架」一致。正文照录原书。]

这样的流形 $M^{4k}$ 的 $A$-亏格由下式给出：

\begin{align}
A_j(M^{4k}) &= 0\quad \text{当 }j<k, \nonumber\\
A_k(M^{4k}) &= a_k\frac{\alpha_k}{\beta_k}\lambda_M(2k-1)!, \tag{3.28}
\end{align}

于是由 $A$-亏格的整性，数 $a_k\dfrac{\alpha_k}{\beta_k}\lambda_M(2k-1)!$ 当 $k$ 为奇数时必须是偶数。由此可知，$|J\pi_{4k-1}(SO)|$ 被分数 $B_k/4k$ 的分母整除，其中 $B_k$ 是第 $k$ 个 Bernoulli 数，因为

\[ \frac{\alpha_k}{\beta_k}=\frac{B_k}{2\cdot(2k)!},\qquad A_k(M^{4k})=\frac{\alpha_k}{\beta_k}p_k+\cdots. \]

（这一结果由 Milnor 与 Kervaire 于 20 世纪 50 年代末证明。）

\begin{example}{a)}
$k=1$。此时 $A_1=-\dfrac{p_1}{24}$，$a_1=2$。由于数 $\lambda_{\min}\cdot\dfrac{2}{24}$ 是偶数，可知 $|J\pi_3(SO)|$ 被 $24$ 整除。（回忆事实上 $\pi_3(SO)\cong\Z/24$。）又因 $A_1(M^4)$ 是偶数且 $\tau=p_1/3$，所以符号差必须被 16 整除（Rohlin，20 世纪 50 年代初）。
\end{example}

\begin{remark}{}
把相交指数看作二次型作代数讨论，特别是利用它的偶值与幺模性，只能得到符号差被 8 整除。\hfill$\square$
\end{remark}

\begin{example}{b)}
$k=2$。此时 $a_2=1$，且

\[ A_2=\frac{1}{45\cdot 2^7}\bigl(7p_1^2-4p_2\bigr), \]

所以 $\lambda_{\min}\cdot\dfrac{1}{240}$ 在此情形必须是整数。因此 $|J\pi_7(SO)|$ 被 240 整除。
\end{example}

% [原文 p.223]
\begin{remark}{}
稳定同伦群 $\pi_{n+j}(S^n)$ 已用 Cartan--Serre--Adams 型方法算出，其范围显著超出 $j=7$，这归功于若干作者的努力（特别是 Toda）。关于 $J\pi_{4k-1}(SO)$ 的阶的公式 (3.27) 对一切 $k$ 成立，这是流形论中的一个重要结果；特别地，它给出球面稳定同伦群阶的一个下界。这种计算球面稳定同伦群的「几何」途径，只有在广义（超常）上同调理论的框架内才能与代数方法（Cartan--Serre 方法与 Adams 谱序列）结合。\hfill$\square$
\end{remark}

我们再提一个符号差在流形论中的应用，它具有根本的重要性。借助符号差公式，人们发现了流形上的非平凡光滑结构，最早是在球面 $S^7$ 上（由 Milnor 于 20 世纪 50 年代末发现）。构造如下。

考虑球面 $B=S^4$ 上一个 Hopf 型的 4 维向量丛 $\eta$。相应的球丛 $S(\eta)$ 的纤维是 $S^3$，结构群是 $SO_4$。设

\begin{equation}
\chi_4(\eta)=\mu,\qquad \langle\mu,[S^4]\rangle=1, \tag{3.29}
\end{equation}

其中 $\mu$ 是 $H^4(S^4;\Z)$ 的生成元。这个条件蕴含纤维为 $S^3$ 的丛的全空间同伦等价于 $S^7$；我们把该全空间记为 $M_\alpha^7$，其中 $\alpha$ 由 $\mu$ 与 $\eta$ 的第一 Pontryagin 类定义为

\[ p_1(\eta)=\alpha\mu,\qquad \langle\mu,[S^4]\rangle=1. \]

以 $N_\alpha^8$ 记相应的、纤维为 $D^4$ 的圆盘丛 $D(\eta)$ 的全空间。注意 $\partial N_\alpha^8=M_\alpha^7$。不难证明 $\alpha$ 必须是偶数：$\alpha=2k$。球丛为标准 Hopf 纤维化 $S^7\to S^4$ 的四元数 Hopf 丛对应于 $k=1$ 的情形。Milnor 对一切偶数 $\alpha$ 显式构造了这些丛，同时在每个这样的流形 $M_\alpha^7$ 上定义了一个恰有两个临界点（一个极大点与一个极小点）的 Morse 函数。由此可知，每个 $M_\alpha^7$ 都同胚（甚至分片线性同胚）于标准球面 $S^7$。现在考虑 PL 流形

\[ \bar N_\alpha^8=N_\alpha^8\cup_{S^7}D^8. \]

假定 $\bar N_\alpha^8$ 是光滑流形，我们来计算它的第一 Pontryagin 类与符号差：

\begin{equation}
p_1(\bar N_\alpha^8)=p_1(N_\alpha^8)=\alpha=2k;\qquad \tau(\bar N_\alpha^8)=\tau(N_\alpha^8)=1. \tag{3.30}
\end{equation}

于是由 Hirzebruch 公式，

\[ p_2(\bar N_\alpha^8)=\frac{45\tau+p_1^2}{7}=\frac{45\tau+(2k)^2}{7}, \]

% [原文 p.224]
当 $k=2$（即 $\alpha=4$）时这不是整数。由此断定流形 $\bar N_4^8$ 不是光滑流形，从而 $M_4^7$ 不微分同胚于 $S^7$。

由 Pontryagin 数的拓扑不变性可知，事实上流形 $\bar N_4^8$ 不同胚于任何光滑流形；但显然 $\bar N_4^8$ 是 PL 流形。关于球面与其他流形上光滑结构理论的进一步发展，将在下一节讨论。

% [Tutor Note: 这里用的是「Pontryagin 数的拓扑不变性」——这是 Novikov 的定理，本书 §5 才正式讨论。正因为 $p_2$ 是拓扑不变量，而整数性又是光滑流形必满足的条件，才能把「$p_2$ 不是整数」提升为「不光滑」而非仅仅「这个光滑结构不对」。]

现在我们回到配边论以及用子流形实现闭链的问题。我们已经看到，通过 Thom--Pontryagin 构造计算配边环，归结为研究万有向量 $G$-丛的 Thom 空间的同伦性质（对 $G=O,\,SO,\,U,\,SU,\,Sp,\,Spin$ 等）。另一方面，Adams 的方法则要计算上同调环 $H^*(M(G);\Z/p)$ 以及 Steenrod 运算在其上的作用，即该上同调环作为 Steenrod 代数上的模的结构（见第 3 章 §6）。除 $H^*(M(\Spin);\Z/2)$ 的情形外，这些模的结构都可以由下述观察推出：底空间作为零截面嵌入 $i:BG_n\to M(G_n)$ 诱导一个（Steenrod 代数上的模的）单射

\[ i^*:H^*(M(G_n);\Z/p)\to H^*(BG_n;\Z/p), \]

其像是由 $i^*(u_n)$ 生成的主理想，其中 $u_n$ 是 $H^*(M(G_n);\Z/p)$ 中的 Thom 类。特别地，这表明乘积算子

\begin{equation}
z\mapsto i^*(u_n)\cdot z=i^*(z\cdot\varphi(1))=i^*\varphi(z) \tag{3.31}
\end{equation}

在 $M(G_n)$ 的相应上同调群中具有零核。

对 $G_n=O_n,\,SO_n$，得到

\begin{align}
i^*(u_n) &= w_n, \nonumber\\
H^*(M(G_n);\Z/2) &\cong w_n\cdot H^*(BG_n;\Z/2). \tag{3.32}
\end{align}

注意：若 $G_n=U_n,\,SU_n$，则 $n=2k$；若 $G_n=Sp_n$，则 $n=4k$。有下列同构：

\begin{align}
H^*(M(SO_{2k});\Z/p) &\cong \chi_{2k}\cdot H^*(BSO_{2k};\Z/p), && p>2, \nonumber\\
H^*(M(U_{2k});\Z/p) &\cong c_k\cdot H^*(BU_{2k};\Z/p), && p\ge 2, \nonumber\\
H^*(M(SU_{2k});\Z/p) &\cong c_k\cdot H^*(BSU_{2k};\Z/p), && p>2, \tag{3.33}\\
H^*(M(Sp_{4k});\Z/p) &\cong i^*(u_{4k})\cdot H^*(BSp_{4k};\Z/p), && p>2. \nonumber
\end{align}

回忆（见上文 §1）$BG_n$ 的上同调环可用符号生成元描述（即 $H^*(T^n)$ 的生成元中的对称多项式，其中 $T^n$ 是 $G_n$ 的极大环面：对 $G_n=O_n,\,SO_n$ 是模 2 的 $t_1,\dots,t_n$，用它们表出 $w_n=t_1\cdots t_n$；其余情形是 2 维生成元 $u_1,\dots,u_k$，用它们表出 $\chi_{2k}=c_k=u_1\cdots u_k$）。

% [原文 p.225]
同构 (3.32) 蕴含：在维数 $\le 2n$ 时，上同调环 $H^*(M(O_n);\Z/2)$ 是 Steenrod 代数 $\cA_2$ 上的自由模，以非二进单项式

\[
\begin{aligned}
\mathfrak{p}&=\sum t_{i_1}^{q_1}\cdots t_{i_s}^{q_s},\\
q&=\sum q_j<n,\qquad q_j\ \text{不形如}\ 2^r-1,
\end{aligned}
\]
为生成元，而这些生成元在 Thom 类 $w_n=t_1\cdots t_n$ 上的作用就是多项式的普通乘法：

\[ \mathfrak{p}(w_n)=\mathfrak{p}\,t_1\cdots t_n\in H^{n+q}(M(O_n);\Z/2),\qquad q=\sum q_j. \]

（这是 Thom 的一个定理。）这一结果允许（借助 Adams 谱序列）完全计算 Thom 空间 $M(O_n)$ 的同伦群（在维数 $\le 2n$ 中），等价地，完全计算谱 $MO$ 的所有稳定同伦群；由 Thom--Pontryagin 构造，后者可等同于配边环 $\Omega_*^{O}=\mathfrak{N}_*$。作为这一计算的一个直接应用，有下述定理：

\begin{theorem}{}
流形 $M^n$ 的（模 2）同调中任何维数 $\le n/2$ 的闭链都可由闭子流形实现；而在所有维数下，闭链都可实现为某个闭流形的像。
\end{theorem}

对 $p>2$ 与 $G=SO,\,U,\,SU,\,Sp$，$\cA_p$-模 $H^*(M(G);\Z/p)$ 是若干个一维 $\cA_p$-模 $\cA_p/(\beta)$ 的直和，其中 $\beta$ 是 Bockstein 算子，$(\beta)$ 是由 $\beta$ 生成的双边理想。模 $\cA_p/(\beta)$ 由关系 $\beta(z)\equiv 0$ 给出。

\begin{remark}{}
后一断言对 $H^*(M(U);\Z/2)$ 也成立，只须把 Bockstein 算子 $\beta$ 换成运算 $Sq^1$。\hfill$\square$
\end{remark}

对 $G=SO,\,SU,\,Sp$ 与 $\Spin$，模 $H^*(M(G);\Z/2)$（作为 $\cA_2$-模）的结构略有不同。它们同样是若干一维 $\cA_2$-模的直和，但被加项有变化：对 $G=SO$，它们是自由的，或者带单个关系 $Sq^1(v)=0$（$v$ 是生成元）；对 $G=SU$，除同样的关系 $Sq^1(v)=0$ 外，还对每个元素 $z$ 都有恒等关系 $Sq^2(z)\equiv 0$；对 $G=Sp$（最难的情形），有两个恒等关系 $Sq^1(z)\equiv 0$、$Sq^2(z)\equiv 0$（对一切 $z$）；最后，对 $G=\Spin$，被加项带有两个关系 $Sq^1(v)=0$、$Sq^2(v)=0$（$v$ 是该 $\cA_2$-模的生成元）。（这些结果属于 Milnor 与 Novikov（1960 年前后），$\Spin$ 的情形属于 Anderson、Brown 与 Peterson，20 世纪 60 年代中期。）

应当指出，上述所有配边环上都有乘积结构（由流形的直积诱导）。这一「几何的」

% [原文 p.226]
乘积结构经由对角同态

\[ \Delta:M\to M\otimes_{\Z/p}M \]

决定了 $\cA_p$-模 $M=H^*(M(G);\Z/p)$ 上的一个代数「乘积」，这又决定了上同调群

\[ \Ext_{\cA_p}^{*,*}(M,\Z/p) \]

上的一个乘法。

在克服若干代数障碍之后，Adams 谱序列（见第 3 章 §6、§7）给出配边环 $\Omega_*^{SO}$、$\Omega_*^{U}$ 的结构（Milnor、Novikov）。应当提到，群 $\Omega_*^{SO}$、$\Omega_*^{SU}$ 的 2-挠早已用不同的、纯几何的想法算出（$\Omega_*^{SO}$ 的情形由 Rohlin 与 Wall 于 20 世纪 50 年代末完成，$\Omega_*^{SU}$ 的情形由 Conner 与 Floyd 于 20 世纪 60 年代中期完成）；然而，把这些几何方法与同调代数的方法结合起来，只有在广义（超常）上同调（配边）理论 $\Omega_U^*(\cdot)$ 与 $\Omega_O^*(\cdot)$ 的框架内才可行。

% [Tutor Note: 原书 p.226 印作「yields up the structures」，此处多一 up，属排印衍字。]

结果如下：任何配边环 $\Omega_*^{G}$（$G$ 如前）都没有 $p>2$ 的 $p$-挠，$\Omega_*^{U}$ 的情形连 $p\ge 2$ 也没有。环 $\Omega_*^{G}\otimes\Z/p$ 与 $\Omega_*^{G}\otimes\Q$ 对 $p>2$ 是多项式环。环 $\Omega_*^{U}$（在整数上）是偶维生成元上的多项式环，每个偶维数有一个生成元。在这些生成元 $[M^{2n}]\in\Omega_{2n}^{U}$ 上，Newton 多项式 $\tilde c_n=n!\,\ch_n$ 取下列值：

\[ \langle\tilde c_n,[M^{2n}]\rangle=\begin{cases}1, & n\ne p^i-1,\\ p, & n=p^i-1,\end{cases}\qquad p\ \text{为素数},\quad [M^{2n}]\in\Omega^U_{2n}. \]

由于 $\langle\tilde c_n,\CP^n\rangle=n+1$，可知只有元素 $[\CP^{p-1}]$ 满足上述条件。所有其他生成元 $[M^{2n}]$ 都可由 Milnor 流形 $H_{r,t}$ 构造出来，它们在乘积 $\CP^r\times\CP^t$ 中代表闭链 $\CP^r\times\CP^{t-1}+\CP^{r-1}\times\CP^t$。元素 $[H_{r,t}]$ 的线性组合给出环 $\Omega_*^{U}$ 的一族完全的乘性生成元。

下述定理（Novikov，20 世纪 60 年代初）是关于稳定维数 $j<n/2$ 中整闭链 $z\in H_j(M^n;\Z)$ 可由定向子流形表示（或者，在任意维数下，可由流形的连续像表示）的定理，它是关于复形 $M(SO_m)$（或谱 $MSO$）的稳定同伦结构的一些结果的推论：

\begin{theorem*}{}
存在奇整数 $\lambda$，使闭链 $\lambda z$ 可以这样表示；此外，若对一切 $k\ge 1$、$p>2$，同调群 $H_j(M^n;\Z)$ 在维数 $j=2k(p-1)-1$ 中没有 $p$-挠，那么事实上 $\lambda$ 可以取为 1。
\end{theorem*}

\begin{remark}{}
若在上述群中有 $p$-挠，则因子 $\lambda$ 可以取为乘积

\[ \lambda=\prod_{p>2}p^{\bigl[\frac{j}{2(p-1)}\bigr]}, \]

此即 Hirzebruch 多项式 $L_m(p_1,\dots,p_m)$ 中一般系数的分母公式，其中 $m=\bigl[\dfrac{j}{4}\bigr]$。\hfill$\square$
\end{remark}

% CHUNK ch04g DONE
